\section{Order, Distance, and Length}
\label{sec:order-distance-length}

\begin{topicbox}{Distance on the Line}
\begin{exposition}
The absolute value of a difference turns $\mathbb{R}$ into a metric space. Length of an interval is the distance between its endpoints. These are the quantitative tools underneath every later definition of nearness.
\end{exposition}
\end{topicbox}
\begin{tcolorbox}[colback=defbox, colframe=defborder, arc=2pt,
  left=6pt, right=6pt, top=4pt, bottom=4pt,
  title={\small\textbf{Definition (Distance on the Real Line)}},
  fonttitle=\small\bfseries]
\begin{definition}[Distance on the Real Line]
\label{def:distance-on-real-line}
The \emph{distance} on $\mathbb{R}$ is the map $d:\mathbb{R}\times\mathbb{R}\to\mathbb{R}$ given by
\[
d(x,y)=|x-y|.
\]
\end{definition}
\end{tcolorbox}
\begin{remark*}[Standard quantified statement]
\[
\forall x,y\in\mathbb{R}\;(d(x,y)=|x-y|).
\]
\end{remark*}
\begin{remark*}[Interpretation]
The distance between two points is the length of the segment joining them. It is the metric that makes ``close'' precise and is the basis for balls, neighborhoods, and open sets.
\end{remark*}
\begin{dependencies}
\begin{itemize}
  \item TODO: link to the definition of absolute value on $\mathbb{R}$.
\end{itemize}
\end{dependencies}
\begin{tcolorbox}[colback=defbox, colframe=defborder, arc=2pt,
  left=6pt, right=6pt, top=4pt, bottom=4pt,
  title={\small\textbf{Definition (Length of a Bounded Interval)}},
  fonttitle=\small\bfseries]
\begin{definition}[Length of a Bounded Interval]
\label{def:length-of-interval}
Let $a,b\in\mathbb{R}$ with $a\le b$. The \emph{length} of any bounded interval with endpoints $a$ and $b$ is
\[
\ell=b-a.
\]
\end{definition}
\end{tcolorbox}
\begin{remark*}[Interpretation]
Length depends only on the endpoints, not on whether they are included. It is the distance $d(a,b)$ and measures the size of the basic pieces of the line.
\end{remark*}
\begin{dependencies}
\begin{itemize}
  \item \hyperref[def:distance-on-real-line]{Distance on the Real Line}
\end{itemize}
\end{dependencies}
\begin{tcolorbox}[colback=thmbox, colframe=thmborder, arc=2pt,
  left=6pt, right=6pt, top=4pt, bottom=4pt,
  title={\small\textbf{Theorem (The Distance Function Is a Metric)}},
  fonttitle=\small\bfseries]
\begin{theorem}[The Distance Function Is a Metric]
\label{thm:distance-is-a-metric}
The map $d(x,y)=|x-y|$ satisfies, for all $x,y,z\in\mathbb{R}$: (i) $d(x,y)\ge 0$ with $d(x,y)=0$ iff $x=y$; (ii) $d(x,y)=d(y,x)$; and (iii) $d(x,z)\le d(x,y)+d(y,z)$. Hence $(\mathbb{R},d)$ is a metric space.
\hyperref[prf:distance-is-a-metric]{Go to proof.}
\end{theorem}
\end{tcolorbox}
\begin{remark*}[Standard quantified statement]
\[
\forall x,y,z\in\mathbb{R}\;\bigl(d(x,y)\ge 0\ \land\ (d(x,y)=0\Leftrightarrow x=y)\ \land\ d(x,y)=d(y,x)\ \land\ d(x,z)\le d(x,y)+d(y,z)\bigr).
\]
\end{remark*}
\begin{remark*}[Interpretation]
The three metric axioms are exactly the properties of $|\cdot|$: nonnegativity with the zero-distance characterization, symmetry, and the triangle inequality. This is the precise sense in which $\mathbb{R}$ is a one-dimensional metric space.
\end{remark*}
\begin{dependencies}
\begin{itemize}
  \item \hyperref[def:distance-on-real-line]{Distance on the Real Line}
  \item TODO: link to the triangle inequality for absolute value.
\end{itemize}
\end{dependencies}
